\section{Tate coordinates and integral frames}\label{h3:sec:towers}

Our goal is to describe the two completed height strata by extension
spaces carrying their actual integral frames.  The finite torsion
coordinates first identify the completed Tate-basis towers.  Relative
sections on the Fargues--Fontaine curve then identify their quotient
frames and determinant characters.  We finish by descending the
markings and the finite root torsors to the algebraic coefficient rings
needed for ordinary descent.

Fix a finite field $k$ containing $\F_{p^6}$.  For $s=1,2,3$ let
$\Gamma_s$ be a Honda formal group over $k$, with
$[p]_{\Gamma_s}(T)=T^{p^s}$, and put
\[
 D_s=\hthreeEnd_k(\Gamma_s)[1/p],\qquad
 P_s=\Aut_k(\Gamma_s)=\mathcal O_{D_s}^{\times}.
\]
We use the convention in which the corresponding Fargues--Fontaine bundle
$V_s=\mathcal O(1/s)$ has rank $s$ and degree one.  The notation $H_s$
denotes the perfected open Honda ball: on an affinoid perfectoid
$k$-algebra $(R,R^+)$ it is $R^{\circ\circ}$, with addition defined by
$\Gamma_s$.  It is a sheaf of $\Q_p$-vector spaces, since $[p]$ is
invertible.  All assertions about a locus in these sheaves are understood
fiberwise at geometric points.

Let $\mathcal G$ be the $p$-divisible group obtained from the universal
height-three deformation, reduced modulo $p$, over $A=k[[x,y]]$.
Here and below this is the algebraic $p$-divisible group formed by its
finite kernels.  In detail, preparation applied to $[p^l](T)$ gives
finite free algebras
\[
 A[[T]]/([p^l](T))
\]
of rank $p^{3l}$.  Their coordinates are topologically nilpotent for the
maximal-ideal topology, so substitution in the formal group law defines
their algebraic group operations.  Preparation applied to
$[p](T)-Z$ proves the required finite flat transition maps.  We base
change these finite groups when we pass to a height stratum; we do not
complete the generic fiber at its identity and thereby discard its
\'etale part.

For $m=1,2$ put $r=3-m$ and
\[
 A_m=A/(u_1,\ldots,u_{m-1})=k[[u_m,\ldots,u_2]],
 \qquad u_1=x,\quad u_2=y,
\]
where the ideal is zero for $m=1$.  We use the same symbol
$\mathcal G$ for its base change to $\operatorname{Spec}(A_m)$ and
for subsequent base changes to covers.  Let $\mathcal X_m$ be the completed perfection of the analytic locus
where the displayed parameters are topologically nilpotent and $u_m$
is invertible.  The $p$-divisible group on this locus has connected
height $m$ and \'etale height $r$.  Let $\mathcal T_m\to\mathcal X_m$
be the tower of integral bases of its \'etale Tate module.  Thus its
structure group is $\GL_r(\Z_p)$.

\subsection{Finite levels and their norms}

\begin{lemma}\label{h3:geom:regular-levels}
For $l\geq1$, form the reduced closure of the locus of $r$ points of
$\mathcal G[p^l]$ whose images are a basis of the \'etale quotient at
the generic point of $A_m$.  Its coordinate algebra $B_l$ is finite
over $A_m$, and its lifted point coordinates $t_{1,l},\ldots,t_{r,l}$
give isomorphisms
\[
 B_l=k[[t_{1,l},\ldots,t_{r,l}]],\qquad
 C_l:=A_m[t_{1,l}^{p^{ml}},\ldots,t_{r,l}^{p^{ml}}]
     =k[[t_{1,l}^{p^{ml}},\ldots,t_{r,l}^{p^{ml}}]].
\]
The subalgebra on the left is already complete.  The algebra
$C_l[u_m^{-1}]$ is the finite \'etale algebra of bases of
$\mathcal G^{\mathrm{et}}[p^l]$ over $A_m[u_m^{-1}]$.
Over this basis cover, $B_l[u_m^{-1}]$ is the algebra of all
generator lifts, without taking a reduction.
The transition maps $B_j\to B_l$, $j\leq l$, are finite injective maps,
and
\[
 t_{i,j}\in
 k[[t_{1,l}^{p^{m(l-j)}},\ldots,t_{r,l}^{p^{m(l-j)}}]].
\]
\end{lemma}

\begin{proof}
One precise definition of the closure is the image of the finite
$A_m$-algebra of $r$-tuples in $\mathcal G[p^l]$ in the reduced
coordinate algebra of the indicated generic locus.  It is therefore
finite and reduced, the map $A_m\to B_l$ is injective, and every
irreducible component dominates $\Spec A_m$.  At the closed point all
torsion coordinates are nilpotent.  Consequently $B_l$ is local with
residue field $k$, and its maximal ideal is generated by the height
parameters and the $t_{i,l}$.

Write $d=p^m$ and $q=p^3$.  Put
$a_i=[p^{l-1}]_{\mathcal G}(t_{i,l})$, and for
$v=(v_1,\ldots,v_r)\in\F_p^r$ put
$a_v=\sum_{\mathcal G}[v_i]a_i$.  If $W_p(T)$ is the distinguished
polynomial for $[p]_{\mathcal G}(T)$, then
\begin{equation}\label{h3:geom:basis-divisor}
 W_p(T)=\prod_{v\in\F_p^r}(T-a_v)^d
       =\prod_{v\in\F_p^r}(T^d-a_v^d).
\end{equation}
Indeed, over each geometric generic field the connected part has
length $d$ and the $a_v$ enumerate the distinct \'etale points.
Both sides are monic of degree $q$.  Equality of their coefficients
there implies equality in the reduced generic closure.

On $A_m$, the series $[p](T)$ factors through $T^d$.
It follows that every $a_v^d$ is a series over $A_m$ in
$w_i=t_{i,l}^{d^l}$, with zero constant term as a series in the $w_i$.
The algebra $C_l=A_m[w_1,\ldots,w_r]$ is finite over the complete ring
$A_m$, hence complete.  Equation~\eqref{h3:geom:basis-divisor} is an
identity over $C_l$.  Modulo $(w_1,\ldots,w_r)$ it becomes $W_p(T)=T^q$.
The preparation unit for $[p](T)$ lies over $A_m$; thus every
coefficient below degree $q$ in $[p](T)$ vanishes in this quotient.
The successive height-coordinate congruences
\[
 [p](T)\equiv u_iT^{p^i}
 \pmod{(u_m,\ldots,u_{i-1},T^{p^i+1})},\qquad m\leq i\leq2,
\]
now imply, successively, that all $u_i$ vanish.  Hence the maximal
ideal of $C_l$ is generated by the $r$ elements $w_i$.
The ring has dimension $r$, because it is integral over $A_m$.
It is a regular local ring.  The surjection
$k[[W_1,\ldots,W_r]]\to C_l$, $W_i\mapsto w_i$, is an
isomorphism: a nonzero kernel would lower dimension in the regular
domain on the left.  Applying the identical maximal-ideal argument
with $t_{i,l}$ in place of $w_i$ proves $B_l=k[[t_{1,l},\ldots,t_{r,l}]]$.
This also proves the asserted identification of the inclusion
$C_l\subset B_l$, rather than just abstract regularity of both rings.

Over $u_m\ne0$, relative Frobenius of order $ml$ identifies the
quotient of $\mathcal G[p^l]$ by its connected part with its image
under $T\mapsto T^{d^l}$.  Here is a check that retains the
prepared finite algebra.  Write $[p^l](T)=g_l(T^{d^l})$ and
prepare $g_l(S)=U_l(S)V_l(S)$, with $U_l$ a unit and $V_l$
monic of degree $p^{rl}$.  Its linear coefficient $a_l$ is a
unit multiple of $u_m^{1+d+\cdots+d^{l-1}}$.  Invariant
differentials for this Frobenius-divided homomorphism give
\[
 g_l'(S)=a_l\frac{h_{\mathrm{source}}(S)}
                         {h_{\mathrm{target}}(g_l(S))},
\]
where both $h$ are unit series.  In the finite algebra defined
by $V_l$, this says that $U_lV_l'$ becomes a unit after
inverting $u_m$.  Thus that algebra is \'etale of rank
$p^{rl}$.  The kernel of relative Frobenius has rank $d^l$
and is the connected part.  Taking bases gives exactly
$S_l=C_l[u_m^{-1}]$.

Over $S_l$, the scheme of lifts of its universal basis is
finite locally free of rank $d^{lr}$.  Its coordinate algebra
surjects onto $B_l[u_m^{-1}]$: the latter is generated by the
same lifted coordinates.  But the identified power-series
rings show that $B_l[u_m^{-1}]$ is free over $S_l$ of this
same rank, with the monomials $\prod_i t_{i,l}^{e_i}$,
$0\leq e_i<d^l$, as a basis.  A surjection between finite
locally free modules of the same rank is an isomorphism.
This proves the assertion about the full lift scheme; in
particular the earlier generic reduction has removed no
infinitesimal part on the exact-height locus.

Every \'etale basis at level $j$ lifts to a basis at level $l$ after a
geometric field extension.  The transition map on generic coordinate
algebras is therefore injective.  Its restriction $B_j\to B_l$ is
injective by generic density and finite since $B_l$ is finite over
$A_m$.  Finally
$t_{i,j}=[p^{l-j}]_{\mathcal G}(t_{i,l})$ factors through
$t_{i,l}^{d^{l-j}}$.  Substituting the expressions for the base
parameters in $C_l$ gives the last assertion.
\end{proof}

The finite-level calculation has retained both the \'etale basis and
every infinitesimal lift.  To pass to the completed tower, we now give
explicit inverse coordinates on the universal cover.

\begin{lemma}\label{h3:geom:universal-cover}
Let $(R,R^+)$ be an affinoid perfectoid $k$-algebra and specialize the
parameters of $A_m$ to $R^{\circ\circ}$.  For the resulting formal law
$\mathcal G_a$, there is a natural additive isomorphism
\[
 H_3(R)\ \xrightarrow{\ \Phi_a\ }\
 \varprojlim_{[p]}\mathcal G_a(R^{\circ\circ}).
\]
Writing $P_a=[p]_{\mathcal G_a}$, its coordinates and inverse are
\begin{equation}\label{h3:geom:cover-formulas}
 \Phi_a(z)_j=\lim_{n\to\infty}
 P_a^{\circ n}\bigl(z^{1/p^{3(j+n)}}\bigr),\qquad
 \Phi_a^{-1}((t_j))=\lim_{j\to\infty}t_j^{p^{3j}}.
\end{equation}
These formulas commute with change of coefficients and with changes
of deformation framing.  Over an algebraically closed perfectoid
field, the kernel of the zeroth projection is the integral Tate
module of the \'etale quotient of $\mathcal G_a$.
\end{lemma}

\begin{proof}
Choose a power-multiplicative spectral norm and $\lambda<1$ bounding
the specialized parameters.  All coefficients of
$\mathcal G_a-\Gamma_3$ and of $P_a(T)-T^q$, $q=p^3$, have norm at
most $\lambda$.  Integral formal series are $1$-Lipschitz on the open
unit ball.  Since $P_a$ factors through $T^d$, $d=p^m$, we have
\[
 |P_a^{\circ n}(v)-P_a^{\circ n}(w)|\leq |v-w|^{d^n}.
\]
Consecutive approximations in the first formula differ by at most
$\lambda^{d^n}$.  Their limits exist, are topologically nilpotent,
and satisfy
\[
 |\Phi_a(z)_j-z^{1/q^j}|\leq\lambda,
 \qquad P_a(\Phi_a(z)_{j+1})=\Phi_a(z)_j.
\]
For a compatible sequence $(t_j)$, consecutive terms in the second
formula differ by at most $\lambda^{q^j}$.  Its limit $z$ satisfies
$|z-t_j^{q^j}|\leq\lambda^{q^j}$.  Substitution in the first formula
and the displayed contraction show in both directions that the
formulas are inverse.  This reasoning never requires a common
strict upper bound on all the $|t_j|$.

The coefficients of $\Gamma_3$ are fixed by the $q$th-power map.
The difference between the two addition laws is at most $\lambda$;
after $n$ applications of $P_a$ it is at most $\lambda^{d^n}$.
Taking the limit proves additivity.  If $g$ changes the deformation
framing, its evaluated coordinate change differs from its Honda
reduction by coefficients of norm at most $\lambda$, after increasing
$\lambda$ if necessary.  The same contraction proves
$\Phi_{ga}(gz)=g\Phi_a(z)$.  Uniform convergence on smaller balls
proves continuity and compatibility with maps of perfectoid
algebras.

Over an algebraically closed field the points with $t_0=0$ are exactly
the compatible $p^j$-torsion points.  A finite connected group over a
perfect field has just its identity as a geometric point, so these
points form the Tate module of the \'etale quotient.  The zeroth
projection is surjective: preparation applied to $[p](T)-w$ gives a
monic distinguished polynomial, each of whose roots is small when
$w$ is small.  Successively choosing roots gives a compatible
sequence with prescribed $t_0$.
\end{proof}

Write $(H_3^r)_{\mathrm{ind}}$ for the locus of tuples whose
images on every geometric fiber are $\Q_p$-linearly independent
for the Honda vector-space structure on $H_3$.

\begin{theorem}[Coordinate comparison]\label{h3:thm:coordinate-comparison}
There is a natural $P_3\times\GL_r(\Z_p)$-equivariant isomorphism
\[
 \mathcal T_m\ \simeq\ (H_3^r)_{\mathrm{ind}}.
\]
The map takes a compatible basis of generator lifts to
$z_i=\lim_l t_{i,l}^{p^{3l}}$.  It identifies the relevant completed
function algebras with their spectral norms on exhausted closed
smaller balls.  The assertion holds on the relative perfectoid site
and remains valid after adjoining arbitrary profinite parameters.
\end{theorem}

\begin{proof}
First retain the reduced finite-level closures, including their
higher-height boundary.  Perfection identifies a point of the
\'etale quotient with its unique connected generator lift, so their
perfected inverse tower is the tower from
\Cref{h3:geom:regular-levels}.

We check the complete algebras after extension to an algebraically
closed perfectoid field $C^\flat$; these field extensions form
v-covers, and all constructions are defined over $k$.
Put $q=p^3$, $d=p^m$ and $y_{i,l}=t_{i,l}^{q^l}$.
For $0<\rho<1$ in the value group, let $X_l(\rho)$ be the closed
polydisc $\max_i|y_{i,l}|\leq\rho$.  If $\lambda$ is the maximum
norm of the height parameters at a point of this polydisc, then
\Cref{h3:geom:regular-levels} at level one gives
\[
 \lambda\leq\|t_{\cdot,1}\|^d,\qquad
 \|t_{\cdot,1}\|\leq\max(\rho^{1/q},\lambda).
\]
The second inequality follows by finite telescoping of the normalized
division coordinates.  If $\lambda>\rho^{1/q}$, the two inequalities
would give $0<\lambda\leq\lambda^d<\lambda$.  Therefore
\begin{equation}\label{h3:geom:uniform-radius}
 \lambda\leq\rho^{d/q}.
\end{equation}
For consecutive levels it follows that
\[
 \|y_{\cdot,l+1}-y_{\cdot,l}\|
       \leq\lambda^{q^l}
       \leq\rho^{d q^{l-1}}<\rho\qquad(l\geq1).
\]
The finite transition maps take $X_{l+1}(\rho)$ to $X_l(\rho)$.
Conversely every point of $X_l(\rho)$ has a lift, by finiteness,
injectivity and lying-over, and the last strict inequality forces
every lift to remain in $X_{l+1}(\rho)$.  They are thus surjective
maps of these closed polydiscs.  Pullback is isometric for their
spectral norms, before and after completed perfection.

In the completed direct limit of the perfected Banach algebras, the
$y_{i,l}$ converge to elements $z_i$.  The homomorphism from the
perfected Gauss algebra of a polydisc of radius $\rho$ sending its
coordinates to $z_i$ is isometric.  Indeed, for any polynomial in
finitely many fractional powers, evaluation on $y_{\cdot,l}$ has
exactly its Gauss norm: $t_{\cdot,l}$ are free polydisc coordinates
and $q^l$ is a power of $p$.  These evaluations converge to evaluation
on $z$.  A nonzero limiting polynomial therefore has precisely that
norm.  Completing preserves the isometry.

Its image is all the completed direct limit.  For $l\geq j$, write
\[
 t_{i,j}=F_{i,j,l}
 (t_{1,l}^{d^{l-j}},\ldots,t_{r,l}^{d^{l-j}}),
 \qquad F_{i,j,l}\in k[[T_1,\ldots,T_r]].
\]
The formulas in \Cref{h3:geom:universal-cover} show that replacing
$t_{i,l}$ by $z_i^{1/q^l}$ changes this expression by at most
\[
 \lambda^{d^{l-j}}
 \leq \rho^{(d/q)d^{l-j}},
\]
which tends to zero.  Its substituted series converges on the
$z$-polydisc.  Thus the closed image contains every $t_{i,j}$ and,
by applying roots to the same estimates, all its $p$-power roots.
These elements generate the completed direct limit.  Varying
$\rho$ proves the comparison for the open balls.

It remains to identify the open locus.  For
$a=(a_1,\ldots,a_r)\in\Z_p^r$, form the compatible points
$t_{a,j}=\sum_{\mathcal G}[a_i]t_{i,j}$.
The normalized limit is $\sum_{\Gamma_3}[a_i]z_i$ by
\Cref{h3:geom:universal-cover}.  If it vanishes, then
$|t_{a,l}|\leq\lambda$ at every level, and hence
\[
 |t_{a,j}|=|[p^{l-j}](t_{a,l})|
       \leq\lambda^{d^{l-j}}\longrightarrow0.
\]
Thus every $t_{a,j}$ vanishes.  Conversely a relation in the Tate
module gives the same relation among the $z_i$.
At height exactly $m$, \eqref{h3:geom:basis-divisor} gives a basis of
the \'etale quotient modulo $p$, and the compatible points give an
injective map $\Z_p^r$ into its Tate module.  At greater height that
Tate module has rank less than $r$, so there is a relation.
Multiplying a rational relation by a power of $p$ reduces to the
integral assertion.  Exact height $m$ is therefore exactly the
independent locus.  This identifies the corresponding open
subsheaves; it is not only a bijection on their field-valued points.

Equivariance is part of \Cref{h3:geom:universal-cover}.  All the
estimates are spectral estimates on affinoid perfectoid algebras.
They apply to continuous functions with a profinite parameter by
taking their uniform norm, so the same isomorphisms retain those
parameters.
\end{proof}

\subsection{Sections and extension spaces on the curve}

We recall precisely the curve results we need.  For a perfectoid
test object $S$, write $X_S$ for its relative Fargues--Fontaine
curve; this notation does not posit a morphism of schemes
$X_S\to S$.  Relative sections and relative cohomology below mean
the corresponding sheaves on the perfectoid site.

\begin{lemma}\label{h3:geom:curve-inputs}
There are natural isomorphisms of additive sheaves
\[
 H_s\simeq\bigl(S\longmapsto H^0(X_S,V_s)\bigr).
\]
The automorphism sheaf of $V_s$ is the locally profinite group
$D_s^{\times}$.  A family fiberwise isomorphic to $V_s$ is locally
of that form for the pro-\'etale topology.  If $\tau$ denotes
the morphism of v-sites associated with $X_S$, then
\[
 R\tau_*\mathcal O^a=\underline{\Q_p}^{\,a},\qquad
 R\tau_*E=\tau_*E[0]
\]
for every family $E$ fiberwise isomorphic to a $V_s$.
Relative sections are fully faithful on $\Q_p$-linear sheaf
morphisms between these bundles and slope-zero bundles.
All assertions commute with base change and v-descent.
\end{lemma}

\begin{proof}
For the unramified degree-$s$ extension of $\Q_p$, the
Lubin--Tate universal cover is the section sheaf of $\mathcal O(1)$
on the curve with that coefficient field
\cite[Proposition II.2.2]{FarguesScholze}.  Unramified pushforward gives $V_s$.
Reduction of the universal cover and tilting identify it with the
Honda universal cover; one can check reduction directly by the
contracting formulas of \Cref{h3:geom:universal-cover}, using
$|[p](v)-[p](w)|\leq\max(|p|\,|v-w|,|v-w|^2)$ before reduction.
This proves the first assertion with its addition and endomorphisms.
For the local-form assertion, the vector-bundle equivalence and
pure-module theorem identify a pure family of slope $c/d$ in lowest terms,
after adjoining $\F_{p^d}$, with a $(c,d)$-$\Q_p$ local
system \cite[Theorems 8.7.8 and 8.5.12]{KedlayaLiu}.  Its semilinear
Frobenius supplies the action of the constant cyclic division
algebra \cite[Definition 8.5.7]{KedlayaLiu}.  These local systems
trivialize in the pro-\'etale topology
\cite[Lemma 9.1.11]{KedlayaLiu}; a module basis over that division
algebra gives the specified basic local form.  Its automorphism
sheaf is therefore the locally profinite $D_s^\times$.
The slope-zero case is also
\cite[Corollary 8.7.10]{KedlayaLiu}.
To retain the ground field $k$, observe that every geometric
Honda endomorphism $f(T)=\sum a_iT^i$ commutes with
$[p](T)=T^{p^s}$.  Hence $a_i^{p^s}=a_i$, so all integral
endomorphisms, and then all of $D_s$, are defined over
$\F_{p^s}\subset k$.  The actual Honda bundle over $k$
therefore defines a map from $B\underline{D_s^\times}$ to
the rank-$s$, degree-one basic locus.  After extending $k$
to its algebraic closure, the preceding local-form theorem
and its identification of automorphisms make this an
equivalence.  Equivalences of v-stacks descend along that
cover, so the frame torsors already have descent over $k$.

For the required cohomology, on the curve with unramified
coefficient field of degree $s$, the bundle $\mathcal O(1)$
is the positive geometric twist of the globally \'etale
bundle $\mathcal O$.  Thus its $H^1$ vanishes on every
affinoid perfectoid test object
\cite[Corollary 8.8.7 and Theorem 8.7.13]{KedlayaLiu}.
Unramified finite pushforward gives the same assertion for
the standard $V_s$.  A family of that basic type first
standardizes pro-\'etale locally, so its $H^1$ sheaf vanishes.
For $\mathcal O^a$, its $H^0$ sheaf is
$\underline{\Q_p}^{\,a}$.  An extension of $\mathcal O$
by $\mathcal O^a$ is pointwise of slope zero; the local-system
equivalence makes it split pro-\'etale locally.  This proves
vanishing of its $H^1$ sheaf, without asserting vanishing
over an unrefined affinoid.  Higher cohomology on affinoid
tests is zero by \cite[Theorem 8.7.13]{KedlayaLiu}.

The needed full faithfulness is relative: for any small v-stack
$S$, the functor
\[
 R\tau_*:\operatorname{Perf}(X_S)\longrightarrow D(S_v,\Q_p)
\]
is fully faithful \cite[Corollary 3.11]{AnschuetzLeBrasFourier}.  For the
basic positive and slope-zero bundles just considered, the
preceding vanishing identifies $R\tau_*E$ with its section
sheaf in degree zero.  Taking
degree-zero morphisms therefore identifies bundle maps with
$\Q_p$-linear maps of section sheaves.  This construction is
compatible with every base change on $S$, so it also applies to
relative forms of the basic bundles.
\end{proof}

Put
\[
 N_m=\bigl(S\longmapsto H^1(X_S,V_m^{\vee})\bigr),\qquad
 \mathcal E_m=(N_m^r)_{\mathrm{ind}}.
\]
The extension interpretation uses sheafified relative $H^1$ and
classifies extensions with both ends framed.  Independence in
$N_m^r$ means independence over $\Q_p$ on every geometric fiber.

\begin{lemma}\label{h3:geom:independent-extensions}
Independent sections of $V_3$ define a subbundle
$\mathcal O^r\hookrightarrow V_3$ with quotient fiberwise $V_m$.
Conversely, an extension of $V_m$ by $\mathcal O^r$ has middle
bundle fiberwise $V_3$ if and only if its class lies in
$\mathcal E_m$.
\end{lemma}

\begin{proof}
First work on an absolute curve.  Let $M$ be the saturation of the
generic span of $r$ independent sections and put $b=\hthreerank M$.
A wedge of $b$ generically independent sections gives a nonzero
section of $\det M$, so $\deg M\geq0$.  Stability of $V_3$ gives
$\deg M\leq b/3<1$.  Since degree is integral, it is zero.
The wedge has no zeros, because a nonzero effective divisor has
positive degree.  The selected sections trivialize $M$.
Since $H^0(X,\mathcal O)=\Q_p$, independence of the original
sections forces $b=r$; hence their map is a subbundle map.
All slopes of its quotient are positive: a quotient of nonpositive
slope would give a nonzero map from $V_3$ to a bundle of slope at
most zero.  A positive bundle of rank $m\leq2$ and degree one is
$V_m$, by the classification of bundles and integrality of degrees
\cite[Theorem II.2.14]{FarguesScholze}.

An extension of $V_m$ by $\mathcal O^r$ has no negative-slope
quotient, since neither end has a nonzero map to a negative bundle.
A nonnegative bundle of rank three and degree one is either $V_3$
or has a quotient $\mathcal O$.  Such a quotient restricts
nontrivially to $\mathcal O^r$, because $\Hom(V_m,\mathcal O)=0$.
The resulting nonzero linear functional on $\Q_p^r$ kills the
extension class.  Conversely, if a nonzero functional kills that
class, its pushout extension splits and gives a quotient
$\mathcal O$.  This proves the equivalence.  Fiberwise
surjectivity of a morphism of vector bundles is a relative
subbundle condition, and \Cref{h3:geom:curve-inputs} supplies local
standard forms and the relative extension sheaf, so the argument
gives the asserted relative loci.
\end{proof}

\subsection{Integral frames and determinant normalization}

The extension locus is now identified as a rational bundle problem.
The next step retains the integral Tate basis and connected marking,
because their determinant lattices determine the group actions in the
coefficient calculation.

A $P_s$-structured form of $V_s$ is a form whose $D_s^\times$-torsor
of frames has been reduced to $P_s$.  Since
\[
 P_s=\hthreeNrd^{-1}(\Z_p^\times),
\]
this is determinant-lattice data.  It is not an integral lattice
in a rank-$s$ slope-zero local system.  We choose determinant
identifications between the $V_s$ over $k$.  They exist there:
the determinant of the cyclic Honda isocrystal differs from the
degree-one rank-one isocrystal by its cyclic permutation sign
$(-1)^{s-1}$.  For even $s$, choose $a\in\F_{p^2}^\times$
with $a^p=-a$; its Teichm\"uller lift changes the Frobenius
multiplier by $-1$.  Thus $\F_{p^2}\subset k$ removes the sign.
We specify a
common valuation normalization in the following proof.

For the height-two statement below, let
$\operatorname{Surj}(V_3,V_2)$ be the relative sheaf of bundle maps
$V_3\to V_2$ that are surjective on every geometric fiber.  Write
$N_2^*=N_2\setminus\{0\}$ for the locus of extension classes that
are nonzero on every geometric fiber.  The zero class represents the
split extension, so this is exactly the locus where the extension of
$V_2$ by $\mathcal O$ is nonsplit on every geometric fiber.

\begin{proposition}[Integral frame comparison]\label{h3:prop:integral-frames}
Put $\mathcal Y_m=[\mathcal X_m/P_3]$.  There is an equivalence
of v-stacks over $k$
\begin{equation}\label{h3:geom:frame-stack}
 \mathcal Y_m\simeq
 [\mathcal E_m/(\GL_r(\Z_p)\times P_m)].
\end{equation}
It carries an actual universal sequence
\begin{equation}\label{h3:geom:universal-sequence}
 0\longrightarrow\mathcal L_r\otimes_{\Z_p}\mathcal O
 \longrightarrow\mathcal V_3
 \longrightarrow\mathcal V_m\longrightarrow0,
\end{equation}
where $\mathcal L_r$ is the integral \'etale Tate local system and
$\mathcal V_s$ are $P_s$-structured forms.  The two presentations
give the classifying maps to $BP_3$, $B\GL_r(\Z_p)$ and $BP_m$,
and their determinant characters satisfy
\begin{equation}\label{h3:geom:determinant-product}
 \hthreeNrd_3=\det_{\GL_r}\,\hthreeNrd_m.
\end{equation}
With the convention $(A,b)e=A_*((b^{-1})^*e)$ for extension classes,
the kernel scalar $u$ acts on $N_m^r$ by $u$ and the quotient scalar
$u$ by $u^{-1}$.

For $m=2$, let $\widetilde{\mathcal X}_2$ be the full
connected-marking tower.  It admits a normalized integral
\'etale generator, fixed by $H=\ker(\hthreeNrd:P_3\to\Z_p^\times)$.
There are compatible equivalences
\begin{equation}\label{h3:geom:normalized-tower}
 \widetilde{\mathcal X}_2\simeq\operatorname{Surj}(V_3,V_2),
 \qquad [\widetilde{\mathcal X}_2/H]\simeq N_2^*.
\end{equation}
The full residual norm quotient acts in the second equivalence by
scalar units, in the kernel convention just specified.  No section
of $P_3\to\Z_p^\times$ is involved.
\end{proposition}

\begin{proof}
Put $R_m=A_m[u_m^{-1}]$.  Apply
\cite[Theorem~1 and its proof, pp.~1--4]{LurieFormalGroups2010}
to the formal completion $\widehat{\mathcal G}_{R_m}$ of $\mathcal G$
over $R_m$ and to $\Gamma_m$, both of exact height $m$.  That theorem defines the
isomorphism ring using every coefficient of the full identity
\[
 F_{\Gamma_m}(f(X),f(Y))=f(F_{\mathcal G}(X,Y))
\]
and expresses it as a filtered union of injective finite \'etale
$R_m$-algebras.  Pullback to $\mathcal X_m$ therefore gives the
pro-\'etale tower of genuine connected markings; any finite set of
Taylor coefficients is defined at one finite stage.  Composition
makes it a torsor under the automorphism sheaf of $\Gamma_m$.
The Honda coefficient calculation in \Cref{h3:geom:curve-inputs}
makes the latter the constant profinite group
$P_m=\Aut_k(\Gamma_m)$: its finite \'etale coefficient stages have
all geometric coefficients in $\F_{p^m}\subset k$.

For the needed bound, pull a genuine marking
$f(T)=\sum_{n\geq1}a_nT^n:\widehat{\mathcal G}_R
\xrightarrow{\sim}\Gamma_{m,R}$ to a characteristic-$p$
affinoid perfectoid test algebra $R$ on this tower, with its
power-multiplicative spectral norm.  Put $q=p^m$.  The factorization
of the $p$-series used in \Cref{h3:geom:regular-levels} reads
\[
 [p]_{\mathcal G}(T)=P(T^q),\qquad
 P(S)=u_mS+\sum_{j\geq2}c_jS^j,
\]
where $u_m\in R^\times\cap R^{\circ\circ}$ and $c_j\in R^\circ$;
these bounds come from the integral deformation coefficients.
The necessary identity $f([p]_{\mathcal G}(T))=f(T)^q$ gives
\[
 a_1^{q-1}=u_m,\qquad
 a_n^q-u_m^na_n=\sum_{1\leq i<n}a_i[S^n]P(S)^i
 \quad(n\geq2).
\]
Here $[S^n]$ denotes coefficient extraction.  First
$\|a_1\|=\|u_m\|^{1/(q-1)}<1$.  If the earlier coefficients have
norm at most one and $A=\|a_n\|>1$, the second equation gives
$A^q\leq\max\{\|u_m\|^nA,1\}\leq A$, a contradiction.
Thus every $a_n$ belongs to $R^\circ$.  The series $f(z)$ converges
in $R^{\circ\circ}$, uniformly on each closed ball of radius less
than one.  The full formal-law identity then evaluates on small
points and proves additivity; evaluation commutes with maps of
perfectoid test algebras.  Compose it with the zeroth projection
in \Cref{h3:geom:universal-cover}.  The resulting additive sheaf map
$H_3\to H_m$ is $\Q_p$-linear: it commutes with integral
formal-group multiplications and with inverse multiplication
by $p$.  It kills the $\Q_p$-span of the Tate basis: it kills
all $p$-power torsion because $[p]$ is injective on $H_m$.
By \Cref{h3:geom:independent-extensions,h3:geom:curve-inputs}, it
therefore factors through the section sheaf of the quotient bundle.
By the relative full faithfulness in \Cref{h3:geom:curve-inputs},
this is the section map of a unique relative bundle morphism
from the quotient to $V_m$.  On every geometric fiber the
section map is nonzero: restrict near the origin, where the
formal isomorphism has nonzero linear coefficient.  After
identifying that fiber of the quotient with $V_m$, the bundle
map is a nonzero element of the division algebra $D_m$, hence
invertible.  A bundle morphism invertible on every geometric
fiber is an isomorphism.  Uniqueness under relative full
faithfulness makes these frames compatible with base change
and with the group actions.

Connected markings form a $P_m$-torsor, and the constructed frames
change by that same maximal compact subgroup.  The relative frame
theorem shows that their only possible discrepancy with the
determinant-lattice reduction of the quotient is its integer
valuation.  This integer is locally constant on the Tate-basis
space and invariant under $P_3$, $P_m$ and $\GL_r(\Z_p)$:
all their determinant changes are units.  For $m=2$, that
space is $H_3\setminus\{0\}$ by
\Cref{h3:thm:coordinate-comparison}.  It is geometrically
connected: after extending the base field it is exhausted by
intersecting connected closed annuli, and perfection preserves
their underlying topology.  The integer is therefore constant,
and we choose the fixed determinant identification to absorb it.
For $m=1$, let the locally constant valuation discrepancy be
$\nu_{\mathrm{val}}$.  Multiplication of the quotient frame by
$p^{-\nu_{\mathrm{val}}}$ corrects it.  Since
$D_1^\times=\Q_p^\times$ is central and $\nu_{\mathrm{val}}$
is invariant under all three actions, this correction is
equivariant and glues without any connectedness assumption.
These choices give the common determinant-lattice normalization.

Now mark both the connected group and the integral Tate basis.
By \Cref{h3:thm:coordinate-comparison} this tower is the sheaf of
exact sequences
\[
 0\to\mathcal O^r\to V_3\to V_m\to0
\]
whose induced determinant identification is a unit multiple of the
chosen one.  The left frame changes are exactly
$\GL_r(\Z_p)\times P_m$.  Forgetting the middle frame instead
gives $\mathcal E_m$, and its middle frames preserving that
determinant lattice form exactly a $P_3$-torsor.  This pair of
torsor presentations proves \eqref{h3:geom:frame-stack}.
Descent of the displayed sequence proves
\eqref{h3:geom:universal-sequence}; taking its determinants proves
\eqref{h3:geom:determinant-product}.  Pushout at the kernel and
pullback at the quotient give the stated action on extension
classes.  This construction uses actual frames throughout, so the
maps to the classifying stacks are part of the equivalence.

For $r=1$, the fixed determinant identification and quotient frame
identify the kernel line with $\mathcal O$.  On the integral
Tate-basis torsor its determinant multiplier is a unit.  Dividing
the generator by that multiplier gives the unique generator with
multiplier one.  Uniqueness glues this normalization and shows
that $g\in P_3$ changes it by $\hthreeNrd(g)$, with the inverse if all
frame conventions are reversed.  It is therefore fixed by $H$.
Forgetting this now determined generator identifies the connected
marking tower with all surjections $V_3\to V_2$.  Indeed, a
surjection $q$ determines its unique kernel generator $i_q$
of determinant multiplier one.  Coordinate comparison on the
entire nonzero Honda ball reconstructs a deformation and a
primitive integral Tate generator from $i_q$, at every valuation.
Scaling $i_q$ by a power of $p$ can change that deformation;
the comparison always returns a basis for the reconstructed one.
Choose a connected marking locally and let $q_h$ be its quotient
frame.  The constant discrepancy already normalized above
ensures that the determinant multiplier of $(i_q,q_h)$ is a
unit.  Write $q=dq_h$ with $d\in D_2^\times$.  Since
$(i_q,q)$ has multiplier one, $v_p(\hthreeNrd d)=0$, so $d\in P_2$.
Thus the required change of marking is an integral Honda
automorphism and yields a unique connected formal isomorphism.
Allowing the middle
bundle to vary while keeping its determinant frame gives precisely
the nonsplit extensions of $V_2$ by $\mathcal O$, namely $N_2^*$.
Changing the middle determinant scales the kernel generator
by its full reduced norm.  This norm maps $P_3$ onto
$\Z_p^\times$, as is seen on the units of the unramified
cubic subfield of $D_3$.
This proves \eqref{h3:geom:normalized-tower}, with the full norm
quotient action.
\end{proof}

\subsection{Algebraic markings and infinitesimal splitting torsors}

The stack comparison supplies normalized markings after completed
perfection.  Ordinary cooperations require those markings at finite
algebraic stages.  We first prove a uniform descent statement that also
allows the number of field factors to grow, and then identify the
resulting finite flat root torsors.

Set $K=k((y))$.  Let $L$ be the algebraic union of finite
connected-marking algebras over $K$, and write
\[
 L^{\mathrm c}=\widehat{L^{\mathrm{perf}}}.
\]
Finite products of complete fields are allowed in this notation;
all assertions below are compatible with their idempotents.  The
commuting actions of $P_3$ and $P_2$ are the actions on the
deformation and on its connected marking respectively.

\begin{lemma}\label{h3:geom:uniform-etale-descent}
Let $L=\colim_i B_i$ be an injective union of finite \'etale
$K$-algebras, with their maximum valuation norms, and put
$C=\widehat{L^{\mathrm{perf}}}$.  For every finite \'etale
$L$-algebra $E$, the map
\[
 \Hom_{L\text{-alg}}(E,L)\longrightarrow
 \Hom_{L\text{-alg}}(E,C)
\]
is bijective.  Each section on the right is defined at one
finite separable stage $B_j$, even if the numbers of field
factors in the $B_i$ are unbounded.
\end{lemma}

\begin{proof}
Every $B_i$ and every $B_i^{1/p^n}$ is a finite product of
complete valued fields.  First let $e\in C$ be an idempotent.
Choose $a\in B_i^{1/p^n}$ with $\|a-e\|<1$.  Then
$\|a\|\leq1$ and $\|a^2-a\|<1$, and $2a-1$ is a unit
whose inverse has norm at most one.  Hensel's lemma in this
one complete finite product gives an idempotent $e_i$ with
$\|e_i-a\|<1$.  Two idempotents at distance less than one
are equal: each of $e(1-e_i)$ and $e_i(1-e)$ is an idempotent
of norm less than one and hence vanishes.  Thus $e=e_i$.
Purely inseparable extensions create no idempotents, so
$e_i\in B_i$.  Any finite list of idempotents descends to
a common stage.

The algebra $E$ descends to a finite \'etale algebra at some
$B_i$.  After a finite idempotent partition it admits a
monogenic \'etale presentation: each component field of
$B_i$ is infinite, so a primitive element may be chosen for
its finite product of separable extensions.  The preceding
paragraph descends the idempotents chosen by a section to
$C$.  It remains to descend its finitely many roots.

Suppose $f(z)=0$ and $f'(z)$ is invertible in $C$, with
$f\in B_i[T]$.  Approximate both $z$ and $f'(z)^{-1}$
simultaneously by $a,c\in B_j^{1/p^n}$, for one $j\geq i$
and one $n$.  We may require $\|1-cf'(a)\|<1$.
The geometric series in this complete finite product makes
$f'(a)$ invertible with inverse norm at most $\|c\|$.
Put $C_0=\max(1,\|c\|)$, and let $M\geq1$ bound the
quadratic Taylor remainder of $f$ on the unit ball about
$a$.  By making the approximations closer, require
$\|f(a)\|C_0^2M<1$ and
$\|f(a)\|C_0<\min(1,(C_0M)^{-1})$.
Newton iteration stays in $B_j^{1/p^n}$, its inverse
derivatives stay bounded by $C_0$, and its residual norms
are bounded successively by
$M C_0^2\|f(a)\|^2$ and their iterates.
It converges there to a root $b$ with
$\|b-a\|\leq C_0\|f(a)\|$.
Taking the original approximation inside the same simple-root
ball makes $b=z$: the divided difference
$(f(b)-f(z))/(b-z)$ is a unit throughout that ball.
The element $b$ is separable over each component of $B_j$
and lies in its purely inseparable extension, so $b\in B_j$.
The finitely many roots use a common $j$.  Injectivity follows
from $L\hookrightarrow C$, proving the assertion.
\end{proof}

\begin{lemma}[Descent of the normalized marking]\label{h3:lem:marking-descent}
The normalized \'etale generator of
\Cref{h3:prop:integral-frames} descends to a compatible integral
marking over $L$.  It is $H$-invariant, and the full $P_3$ action
changes it by the constant unit $\hthreeNrd(g)$, up to the common inverse
convention.  At level $l$, the torsor of its lifts to
$\mathcal G[p^l]$ has coordinate algebra
\begin{equation}\label{h3:geom:height-two-roots}
 L^{1/p^{2l}}.
\end{equation}
The connected marking makes this a torsor for
$\Gamma_2[p^l]_L$.  The group $H$ acts trivially on the translation
group, while $P_2$ acts by constant Honda automorphisms.

For a finite connected-marking level $B=L^\Pi$, with $\Pi\subset P_2$
open, put $B^{\mathrm b}=\widehat{B^{\mathrm{perf}}}$; the superscript
here denotes completed perfection, not tilting.  There is a
pro-\'etale $\Pi$-torsor
\begin{equation}\label{h3:geom:completed-level-torsor}
 \operatorname{Spa}(L^{\mathrm c})\longrightarrow
 \operatorname{Spa}(B^{\mathrm b}),
\end{equation}
and consequently
\begin{equation}\label{h3:geom:finite-level-stack}
 [\operatorname{Spa}(B^{\mathrm b})/H]\simeq[N_2^*/\Pi].
\end{equation}
These equivalences are natural under level changes and after
extension to an algebraically closed perfectoid base.
\end{lemma}

\begin{proof}
For each $l$ the normalized generator is a point of
\[
 \operatorname{Isom}_L
 ((\Z/p^l)_L,\mathcal G^{\mathrm{et}}[p^l]_L).
\]
By \Cref{h3:geom:uniform-etale-descent}, it descends to $L$ at
one finite separable stage.  Compatibility and equivariance
descend through the injective coefficient extension, so the
descended points give the asserted integral marking.  This
argument descends a finite \'etale marking, not a splitting of
the connected--\'etale extension.

The fiber over this marked generator in $\mathcal G[p^l]$ is a
torsor for its connected kernel.  Over the \'etale generator
cover, \Cref{h3:geom:regular-levels} identifies its two fields as
\[
 k((t_l^{p^{2l}}))\subset k((t_l)).
\]
They have purely inseparable degree $p^{2l}$.  The marked
generator supplies a map from the separable field on the left
to $L$.  Root extensions commute with separable scalar
extension, so the torsor algebra after this base change is
exactly \eqref{h3:geom:height-two-roots}; no degree is lost.
The connected marking identifies its group with
$\Gamma_2[p^l]$.  The equivariance already proved fixes both
markings under $H$ and changes them by constant automorphisms
under $P_2$, giving the stated actions on translations.

To prove the torsor assertion, choose a cofinal system of
normal open subgroups $\Pi'\subset\Pi$.  The algebra
$B_{\Pi'}=L^{\Pi'}$ is a finite \'etale $\Pi/\Pi'$-torsor over
$B=L^\Pi$.
Perfection commutes with finite \'etale base change, and
completion commutes with finite modules.  Consequently
\[
 A_{\Pi'}:=B_{\Pi'}\otimes_B B^{\mathrm b}
       \simeq\widehat{B_{\Pi'}^{\mathrm{perf}}}
\]
is again a finite \'etale $\Pi/\Pi'$-torsor.  Their transition
maps are surjective.  The affinoid perfectoid inverse limit
is $\operatorname{Spa}(L^{\mathrm c})$: its algebra is
the uniform completed union of the $A_{\Pi'}$, equal to
$\widehat{L^{\mathrm{perf}}}$.
Passing the finite torsor identities to inverse limits gives
\[
 \operatorname{Spa}(L^{\mathrm c})
 \mathbin{\times}_{\operatorname{Spa}(B^{\mathrm b})}
 \operatorname{Spa}(L^{\mathrm c})
 \simeq\underline\Pi\times\operatorname{Spa}(L^{\mathrm c}).
\]
Each geometric fiber is a nonempty inverse limit of finite
nonempty sets with surjective transition maps.  Thus the
map is surjective and is the pro-\'etale $\Pi$-torsor in
\eqref{h3:geom:completed-level-torsor}.  These arguments permit
arbitrarily many factors at later finite stages.
Finally $H$ and $\Pi$ commute,
so quotienting this torsor and applying
\eqref{h3:geom:normalized-tower} gives
\eqref{h3:geom:finite-level-stack}.  No cohomological Galois-descent
assertion is used in this deduction.
\end{proof}

We fix the periodicity convention needed for the coefficient
arguments.  The invariant cotangent line is
$\omega=\pi_2E_3$; if $|u|=-2$, its generator is $U=u^{-1}$.
In particular $v_1=xU^{p-1}$ modulo $p$.

\begin{lemma}[Invariant differential]\label{h3:lem:invariant-differential}
There is an equivariant isomorphism
\begin{equation}\label{h3:geom:canonical-lines}
 \bigwedge^2\Omega^1_{A/k,\hthreects}\simeq\omega^3[\det],
 \qquad
 \Omega^1_{K/k}\simeq\omega_K^{p+2}[\det].
\end{equation}
Here $\det$ is the reduced norm character of $P_3$.
The connected marking over $L$ supplies a generator $U_0$ of
$\omega_L$ fixed by $P_3$.  Consequently
\eqref{h3:geom:canonical-lines} gives a nonzero differential
$\eta\in\Omega^1_{L/k}$ fixed by $H$, defined at a finite
separable marking level.  Changes of the connected marking
act on these frames by their constant cotangent characters.
\end{lemma}

\begin{proof}
The integral canonical-line theorem gives the first isomorphism
\cite[Theorem 20]{GHdual}.  In that theorem $\omega=E_2$, and
the determinant is the reduced norm, so its line and action
conventions are the ones just fixed.  If $g(U)=q_gU$, comparison
of the first nonzero graded $p$-series coefficient gives
$g(x)=q_g^{-(p-1)}x$ modulo $p$.  Equivalently the graded
coefficient $v_1=xU^{p-1}$ is invariant there.  Hence
$(x)/(x^2)\simeq\omega^{-(p-1)}$ on the height-two slice.
Taking determinants in its conormal sequence yields
\[
 \Omega^1_{k[[y]]/k,\hthreects}
 \simeq\omega^{3+(p-1)}[\det]
 =\omega^{p+2}[\det].
\]
Localizing at $y$ proves the second isomorphism.

There is no change of differential theory at the generic point.
Every power series in one variable over a perfect field can
be grouped uniquely by its exponent modulo $p$; thus $y$ is
a $p$-basis of $K$ and $\Omega^1_{K/k}=K\,dy$.
Separable extension gives
$\Omega^1_{L/k}=L\otimes_K\Omega^1_{K/k}$.
Transport a nonzero invariant differential on $\Gamma_2$
along the tautological connected marking.  This gives $U_0$.
The $P_3$ action transports the marking, so it fixes this
frame.  The image of $U_0^{p+2}\otimes1_{\det}$ under
\eqref{h3:geom:canonical-lines} is $\eta$.  Its remaining
character is the reduced norm, trivial on $H$.  Its
coefficient and that of $U_0$ occur at a finite separable
stage, since each belongs to the algebraic union $L$.
The action of $P_2$ is the constant change of Honda marking,
and differentiation gives the asserted constant characters.
\end{proof}

\begin{lemma}[The ordinary splitting torsor]\label{h3:lem:ordinary-torsor}
Let $B_0=A[x^{-1}]$.  The ordinary connected--\'etale sequence
of $\mathcal G$ over $B_0$ is
\[
 0\longrightarrow\mathcal C\longrightarrow\mathcal G
 \longrightarrow\mathcal E\longrightarrow0,
\]
where $\mathcal C$ has height one and is of multiplicative
type, and $\mathcal E$ is \'etale of height two.
For every $l\geq1$, the scheme of splittings of its
$p^l$-torsion sequence is canonically
\begin{equation}\label{h3:geom:ordinary-root-torsor}
 \Spec B_0^{1/p^l}\longrightarrow\Spec B_0.
\end{equation}
It is a finite flat torsor for
\[
 \operatorname{Hom}(\mathcal E[p^l],\mathcal C[p^l]),
\]
a form of $\mu_{p^l}^2$.  These identifications commute
with the $P_3$ action, with truncation in $l$, and with
Frobenius.  They therefore define a natural tower of
multiplicative-type torsors over the uncompleted ring $B_0$.
\end{lemma}

\begin{proof}
On the ordinary locus the relative Frobenius kernels define
the connected height-one subgroup; dividing by them gives
the finite \'etale quotients.  These compatible finite
groups give $\mathcal C$ and $\mathcal E$.  A height-one
one-dimensional connected $p$-divisible group becomes
$\mu_{p^\infty}$ on its connected-marking cover; at each
finite level the required cover is \'etale.  Thus
$\mathcal C[p^l]$ is of multiplicative type.

The $p^l$-torsion sequence is fpqc exact.  Locally choose
a basis of $\mathcal E[p^l]$ and lift its two generators
to $\mathcal G[p^l]$.  Every lift is killed by $p^l$, so
the two lifts define a homomorphism and a splitting.
Two splittings differ by a unique homomorphism to
$\mathcal C[p^l]$.  Consequently the splitting scheme
is a finite flat torsor for the displayed group, of rank
$p^{2l}$.

The \'etale basis cover is
$B'=C_l[x^{-1}]$, with $m=1$ in
\Cref{h3:geom:regular-levels}.  On this cover, the splitting
scheme is the scheme of generator lifts, whose algebra is
\[
 B_l[x^{-1}]
 =k[[t_{1,l},t_{2,l}]][x^{-1}]
 =(C_l[x^{-1}])^{1/p^l}=(B')^{1/p^l}.
\]
Relative Frobenius on an \'etale algebra is an isomorphism.
Therefore
\[
 (B')^{1/p^l}\simeq
 B'\otimes_{B_0}B_0^{1/p^l}.
\]
These local identifications descend canonically.  More
explicitly, in the splitting algebra the $p^l$th power
of every element lies in the base algebra, as can be
checked after the faithfully flat cover $B'$.  Send
that element to the unique $p^l$th root of its power in
$B_0^{1/p^l}$.  The resulting algebra map is an
isomorphism after the cover, hence an isomorphism.
This description shows both its uniqueness and its
independence of all basis and connected-marking choices.

The splitting construction is functorial for changes of
deformation framing.  Its canonical root identification
is therefore $P_3$-equivariant.  Restricting a splitting
at level $l+1$ to level $l$ is the quotient by the kernel
of the corresponding homomorphism of the translation
groups.  Under the unique root identifications it is
the inclusion $B_0^{1/p^l}\subset B_0^{1/p^{l+1}}$.
Frobenius compatibility follows from the same uniqueness
and the natural relative Frobenius maps of the finite
groups.  Thus all the stated compatibilities hold over
$B_0$ itself, rather than only on an auxiliary \'etale
extension.
\end{proof}
